\section{Extended Language Examples}
\label{app:language}
\subsection{Dialects, Operations, and Attributes}
\label{extended:sec:dialects}
A dialect groups user-defined operations with their operand signatures,
attribute schemas, and semantics. Users define an operation once and import
it into multiple rules; different dialects can provide different operation
vocabularies. Attributes configure an operation independently of its tensor
operands. For example, a transpose carries a permutation list. The following
excerpt declares its signature and attribute in a dialect file,
\code{tensor.tepl}:
\begin{codeblock}[H]
\begin{lstlisting}[language=TEPL]
dialect TensorLang {
  op transpose(input: tensor) -> tensor {
    attrs {
      permutation: index[];
    }
  }
}
\end{lstlisting}
\end{codeblock}
This excerpt introduces the schema; it omits the semantic blocks. In a
semantic definition, \code{dtype} specifies element-type behavior,
\code{axes} declares coordinate bindings, legality, and output order, and
\code{semantics} gives the value at those coordinates. For transpose,
\code{dtype} preserves the element type. Transpose changes output coordinate
order while preserving the coordinate-indexed input value, so a value clause
\code{yield input} does not imply an unchanged tensor interface. We defer
the representation and composition of these blocks to Axes IR in
Section~\ref{sec:axes-ir}.


\subsection{Imports and Rewrite Rules}
A rule imports a dialect and selects the operations it uses:
\begin{codeblock}[H]
\begin{lstlisting}[language=TEPL]
from "tensor.tepl" import TensorLang as t;
use t::{transpose};
\end{lstlisting}
\end{codeblock}
Rewrites are written as algebraic equations with an LHS pattern followed by an
RHS replacement. Nested operations denote producer--consumer edges, while
tensor captures identify matched values or subcomputations. Repeated captures
refer to the same matched value. This application graph is distinct from the
semantic graph used to reason about tensor elements and their coordinates.

Our running pattern is
\code{(transpose[@outer] (transpose[@inner] X))}.
Its intended replacement is a single transpose of \code{X}; the remaining
task is to construct that transpose's permutation. Tensor declarations may
also constrain shape and dtype, but equal dimension sizes alone do not
establish coordinate identity.


\subsection{Attribute Binding and Reuse}
\label{extended:sec:binding}
The annotation \code{op[@name]} binds an operation's attribute descriptor on
the LHS. If LHS attributes are omitted, the matcher captures each operation's
descriptor independently and anonymously. A name is needed when an expression
or an RHS operation refers to that descriptor. Captures are available in
\code{where} conditions, \code{derive} expressions, and RHS operations.
Field access such as \code{@outer.permutation} selects an individual
attribute where the expression context supports it.


Reusing \code{@name} on a compatible RHS operation reuses its descriptor;
a derived descriptor can configure an RHS operation in the same way.
Compatibility checks the attribute structure, rather than the mathematical
correctness of that choice. In the transpose-composition example, binding
\code{@inner} and \code{@outer} makes both permutations available, but copying
either one alone generally does not give the replacement permutation.

Figure~\ref{fig:attribute-binding} illustrates capture and derivation through
a separate transpose-absorption example. The input graph supplies \code{@tr}
and \code{@d}; a user-written \code{f_transform} computes \code{@new_d}.
The table distinguishes the two captured descriptors from the derived one.
Its short fields \code{perm}, \code{l_dim}, and \code{r_dim} belong to a
schematic dialect, rather than the generalized-dot schema used later.
For the displayed permutation \code{[1,0,2]}, the left contraction position
changes from \code{1} in the transposed operand to \code{0} in the original
operand. The figure illustrates descriptor dataflow for this instance; an
arbitrary transpose-absorption rule still requires semantic checking.



\subsection{Writing Attribute Changes}
Conditions and replacement attributes are metadata computations. Expressions
can read captured attributes and available shape and dtype bindings, but they
do not execute tensor computations. Local bindings, builtins, lists, indexing,
conditionals, and comprehensions provide the computations needed to build
descriptors. For transpose composition,
\code{gather(@inner.permutation, @outer.permutation)} computes the replacement
permutation by selecting entries of the inner permutation in outer-permutation
order. Descriptor constructors collect computed fields into immutable
attributes for an RHS operation. Runtime expressibility and semantic
checkability are distinct: some expressions can execute even when the
compiler cannot establish that they preserve a rule's tensor result.


\subsection{Conditions with \texorpdfstring{\code{where}}{where}}
\label{extended:sec:where}
A \code{where} block ends with the condition deciding whether a rule applies.
Any local calculations precede that condition. Only \code{let} bindings end
with a semicolon. Conditions can use captured metadata, builtins, and local
expressions, combining predicates with \code{&&} and \code{||}. For the running
transpose pattern, an illustrative clause is:
\begin{codeblock}[H]
\begin{lstlisting}[language=TEPL]
where {
  let rank = len(@inner.permutation);
  len(@outer.permutation) == rank
}
\end{lstlisting}
\end{codeblock}
The length test illustrates field access and local computation. For a legal
composition of transposes, the equality already follows from operation
legality. A rule can omit \code{where} when no additional condition is needed.

Conditions may express semantic applicability or application policy; TEPL does
not infer profitability decisions. They can also state axis-routing
relationships, for example:
\begin{codeblock}[H]
\begin{lstlisting}[language=TEPL]
@outer.lhs_batching_dimensions == @inner->B
\end{lstlisting}
\end{codeblock}
Its mapping interpretation is introduced in Section~\ref{extended:sec:partitioning}. For a verified semantic rule,
every correctness obligation must follow from operation legality and supported
user-written conditions.


\subsection{Attribute Construction with \texorpdfstring{\code{derive}}{derive}}
The \code{derive} block names replacement descriptors computed from captured
information. Transpose composition combines two captures and an explicit
descriptor constructor:
\begin{codeblock}[H]
\begin{lstlisting}[language=TEPL]
rule compose_transposes_explicit {
  (transpose[@outer] (transpose[@inner] X))
  => (transpose[@combined] X)
  derive {
    @combined = {
      permutation: gather(@inner.permutation, @outer.permutation)
    };
  }
}
\end{lstlisting}
\end{codeblock}
If the inner permutation is $\pi$ and the outer permutation is $\tau$, the resulting
permutation is $\eta[i]=\pi[\tau[i]]$. The three parts of a rule have distinct roles:
\code{where} decides applicability, \code{derive} constructs attributes, and
the RHS specifies the replacement computation.

Local bindings can name intermediate metadata computations. A spread such as
\code{{ ..@outer, permutation: permutation }} copies fields before applying
explicit overrides. The \code{where} and \code{derive} blocks have separate
local scopes; derivations can read captured descriptors and earlier derived
descriptors. With the transpose semantics, this explicit constructor is within
the supported semantic-verification fragment. Checking its attribute types
alone would not establish that the rewrite preserves the result.


\subsection{Automatic Attributes with \texorpdfstring{\code{auto}}{auto}}
\label{extended:sec:auto}
Users can omit replacement attributes when operation semantics support
synthesis. An explicit \code{[auto]} annotation requests the same inference:
\begin{codeblock}[H]
\begin{lstlisting}[language=TEPL]
rule compose_transposes_auto {
  (transpose (transpose X)) => (transpose X)
}
\end{lstlisting}
\end{codeblock}

Matching supplies tensor captures and LHS descriptors. Semantic compilation
then derives a reusable program for the RHS permutation. The explicit rule
specifies how to combine two permutations; the automatic rule asks the compiler
to derive that computation from transpose semantics, again obtaining
$\eta[i]=\pi[\tau[i]]$.

Each omitted RHS descriptor is inferred independently. A named annotation
\code{[@name]} instead reuses an existing descriptor, while
\code{[auto, field = value]} pins one field during inference. Users still
provide operation semantics and any required correctness conditions.
Unsupported inference is rejected with a diagnostic. Constructor synthesis
occurs when the rule is compiled; application evaluates the completed
constructor on the matched metadata. Completed constructors can execute in
the greedy rewrite engine or in generated Rust with the shared runtime.
C++ generation remains unavailable.

Figure~\ref{fig:introducing-auto} makes the compile-time/runtime boundary
explicit using matrix transpose absorption. The compiler synthesizes and
checks a function $f(p)$ from the rule and operation semantics. Each runtime
match supplies the transpose permutation $\pi$ and dot descriptor $\delta$, both contained in $p$, and the
function supplies the replacement dot descriptor. In the matrix instance,
absorbing permutation \code{[1,0]} changes the left contraction position from
$0$ to $1$. The figure shows only that changing field; it does not imply that
other descriptor fields can be ignored in a complete operation definition.



\subsection{From \texorpdfstring{\code{auto}}{auto} to Attribute Synthesis}
The \code{auto} annotation lets users leave RHS attributes to the compiler,
but the compiler must still determine how to construct them. Matching exposes
known attributes and tensor metadata from the LHS. The challenge is to turn
that information into a program that computes the missing RHS attributes
across ranks and shapes and preserves the computation under the user's
conditions. Section~\ref{sec:problem} defines this problem and its correctness contract.


\section{Axes IR Definitions and Examples}
\label{app:axes}
Axes IR is TEPL's intermediate representation of tensor computations. It
describes how operations use tensor axes and compute output values.
Translating operation definitions into this shared representation lets the
compiler track axis relationships across operations, infer replacement
attributes, and check whether completed rewrites preserve the result.


\subsection{Symbolic Axis Groups}
\label{extended:sec:axis-groups}
A symbolic axis group represents an ordered sequence of axes without fixing
its length. One group can, for example, describe all contraction axes of an
operation. This representation borrows arbitrary-rank grouping from
TensorRight's aggregated axes~\cite{tensorright}. TensorRight represents an
aggregated axis as a set of named axes; TEPL uses an ordered sequence and
tracks position mappings from groups to tensor axes for attribute construction.


TEPL declares a group with \code{bind}, associating it with axis positions in
one or more operands. When two operands bind the same group, their selected
axes are paired in list order. Attributes determine those positions, allowing
the same definition to describe different ranks and layouts. The two selections
must have equal length and equal sizes at corresponding positions. Matching
sizes alone do not establish pairing: the binding specifies which coordinates
correspond.

Figure~\ref{fig:axis-groups} illustrates arbitrary-length groups, one concrete
dot binding, and output construction. In that instance, the contraction axes
appear in different physical orders in the two operands. Selecting \code{[2,3]}
in \code{X} and \code{[5,2]} in \code{Y} nevertheless identifies the same
ordered contraction group. Reducing it leaves the ordered output
$B\catseq I\catseq J$.


We use four related notions. A tensor's ordered axes form an interface
$\mathsf{I}=\langle a_0,\ldots,a_{n-1}\rangle$ in dimension order, together with established
mappings from groups into $\mathsf{I}$. Axis labels denote logical coordinates, rather
than sizes. A group $G=\langle g_0,\ldots,g_{m-1}\rangle$ is an ordered
sequence whose order can differ from dimension order. A position sequence
$q=\langle q_0,\ldots,q_{m-1}\rangle$, with $0\leq q_j<n$, locates that
group by $g_j=a_{q_j}$. Positions are zero-based. Finally, $\gath(\mathsf{I},q)$
selects axes in the order specified by $q$, so a mapping records
\begin{equation}
 \gath(\mathsf{I},q)=G.
 \label{extended:eq:mapping}
\end{equation}


The following \code{dot_general} definition declares four attributes and their
axis bindings. Dtype and value semantics are omitted from this excerpt:

\begin{codeblock}[H]
\begin{lstlisting}[language=TEPL]
op dot_general(lhs: tensor, rhs: tensor) -> tensor {
  attrs {
    lhs_batching_dimensions: index[];
    rhs_batching_dimensions: index[];
    lhs_contracting_dimensions: index[];
    rhs_contracting_dimensions: index[];
  }
  axes {
    let lb = attrs.lhs_batching_dimensions;
    let rb = attrs.rhs_batching_dimensions;
    let lc = attrs.lhs_contracting_dimensions;
    let rc = attrs.rhs_contracting_dimensions;
    let lf = ordered_complement(all_axes(lhs), lb ++ lc);
    let rf = ordered_complement(all_axes(rhs), rb ++ rc);
    bind B { lhs: lb, rhs: rb };
    bind K { lhs: lc, rhs: rc };
    bind I { lhs: lf };
    bind J { rhs: rf };
    yield B ++ I ++ J;
  }
}
\end{lstlisting}
\end{codeblock}
Here, \code{all_axes(t)} lists all positions of tensor \code{t}, and \code{++}
concatenates sequences. The builtin \code{ordered_complement} selects remaining
positions in operand order. The axis \code{yield} specifies output order.

As a second concrete instance, smaller than the one in
Figure~\ref{fig:axis-groups}, consider:
\begin{codeblock}[H]
\begin{lstlisting}
lhs:    [i, k0, b, k1]
rhs:    [k1, j, k0, b]
output: [b, i, j]

lhs_batching_dimensions = [2]
rhs_batching_dimensions = [3]
lhs_contracting_dimensions = [1, 3]
rhs_contracting_dimensions = [2, 0]
\end{lstlisting}
\end{codeblock}
The batch binding pairs position $2$ of \code{lhs} with position $3$ of
\code{rhs}. The contraction lists pair $(1,2)$ for $k_0$ and $(3,0)$ for
$k_1$. Remaining axes give $I=\langle i\rangle$ and $J=\langle j\rangle$;
$K$ is reduced and therefore omitted from the output. In mapping notation,
\begin{equation}
\begin{aligned}
\mathsf{I}_{\mathrm{lhs}}&=\langle i,k_0,b,k_1\rangle,
 &q_{\mathrm{lhs}}&=\langle1,3\rangle,\\
\mathsf{I}_{\mathrm{rhs}}&=\langle k_1,j,k_0,b\rangle,
 &q_{\mathrm{rhs}}&=\langle2,0\rangle,\\
G_{\mathrm{lhs}}&=G_{\mathrm{rhs}}=K=\langle k_0,k_1\rangle.
\end{aligned}
\end{equation}
Consequently,
\begin{equation}
 \gath(\mathsf{I}_{\mathrm{lhs}},q_{\mathrm{lhs}})
 =K=\gath(\mathsf{I}_{\mathrm{rhs}},q_{\mathrm{rhs}}).
\end{equation}
Although $K$ has two axes here, neither group length nor position values need
be known at rule-compilation time. Grouping preserves the axes separately and
in order; it does not flatten them into one dimension. Groups can be empty
where an operation permits it, such as an unbatched dot's batch group. Their
identities are scoped so that unrelated coordinates with equal extents remain
distinct. Mappings let the compiler recover where each logical group occurs
and use those positions as replacement attributes.


\subsection{Expressing Tensor Computations}
\label{extended:sec:value-expressions}
The \code{axes} block defines relationships and output order, while
\code{semantics} defines each output element's value. Axes IR distinguishes
\emph{value expressions}, which compute scalar tensor elements, from
\emph{index expressions}, which calculate the indices used to read them.
Ordinary operand references use coordinates established by bindings. Explicit
index expressions handle calculated accesses, including slices and
convolutions. \code{Read} connects the two kinds of expression by reading an
operand at explicitly calculated indices. A \code{semantics} block returns a
value expression with \code{yield}.

For example, elementwise addition without broadcasting binds every axis of
\code{lhs} to the corresponding axis of \code{rhs}, preserving that group:
\begin{codeblock}[H]
\begin{lstlisting}[language=TEPL]
axes {
  bind S {
    lhs: all_axes(lhs),
    rhs: all_axes(rhs)
  };
  yield S;
}
semantics {
  yield Add(lhs, rhs);
}
\end{lstlisting}
\end{codeblock}
The shared binding requires equal rank and matching sizes. At each coordinate
in $S$, the operand references read corresponding elements and \code{Add}
produces their sum. Table~\ref{extended:tab:value-expressions} lists the value primitives.
In this language, \code{Reduce} means summation; it does not imply support
for other reductions such as maximum.

\begin{table}[htbp]
\caption{Value-expression primitives.}
\label{extended:tab:value-expressions}
\begin{tabularx}{\linewidth}{@{}lY@{}}\toprule
Expression & Meaning\\\midrule
\code{input} & Read an operand element using its axis bindings.\\
\code{Const(c)} & Produce a signed integer constant.\\
\code{Add(a, b)} & Add scalar values.\\
\code{Product(a, b)} & Multiply scalar values.\\
\code{Reduce(G, e)} & Sum $e$ over all coordinate combinations of group $G$.\\
\code{Read(input, indices)} & Read an operand at explicitly calculated indices.\\\bottomrule
\end{tabularx}
\end{table}

Index expressions produce an index sequence with one entry for each accessed
axis. Table~\ref{extended:tab:index-expressions} gives their forms. A slice, for example,
uses
\begin{codeblock}[H]
\begin{lstlisting}[language=TEPL]
Read(input,
  IndexOffset(attrs.starts,
    IndexScale(attrs.strides, Coords(O))))
\end{lstlisting}
\end{codeblock}
Here $O$ is the output group. Per axis, the index expression computes
\code{starts + strides * output_indices}. The \code{Read} expression retrieves
the selected value. The index calculation determines the location, and the
enclosing value expression determines the scalar result.

\begin{table}[htbp]
\caption{Index-expression primitives.}
\label{extended:tab:index-expressions}
\begin{tabularx}{\linewidth}{@{}lY@{}}\toprule
Expression & Meaning\\\midrule
\code{Coords(G)} & Current coordinates of $G$ in group order.\\
\code{IndexAdd(a, b)} & Add index sequences elementwise.\\
\code{Concat(a, b)} & Concatenate index sequences.\\
\code{IndexScale(attrs.strides, indices)} & Scale each index by its corresponding entry.\\
\code{IndexOffset(attrs.starts, indices)} & Add the corresponding offset to each index.\\\bottomrule
\end{tabularx}
\end{table}

Value and index expressions share the groups declared in \code{axes}.
Bindings determine participating operand elements; references bound to the
same group use corresponding coordinates. A reduction ranges over its group
while holding the remaining output coordinates fixed. Because a group can
have variable length, one expression can sum over different numbers of
dimensions. The value semantics completing the previous dot definition are:
\begin{codeblock}[H]
\begin{lstlisting}[language=TEPL]
semantics {
  yield Reduce(K, Product(lhs, rhs));
}
\end{lstlisting}
\end{codeblock}
The product reads coordinates determined by $B,I,J,K$, and the reduction sums
over $K$. Output coordinates belong to $B,I,J$, in the declared output order.
The smaller concrete example therefore computes
\begin{equation}
 \operatorname{out}[b,i,j]
 =\sum_{k_0,k_1}\operatorname{lhs}[i,k_0,b,k_1]
                    \operatorname{rhs}[k_1,j,k_0,b].
\end{equation}

These primitives describe elementwise operations, transpose, summation,
and dot products. Slices and supported convolutions calculate input indices;
broadcasting adds output axes while reusing input values along them. Output
shape remains separate from element computation: a tensor filled with zeros,
for example, still needs a specified shape.


\subsection{Reasoning About Composed Operations}
\label{extended:sec:composition}
Value semantics provide a common language for tensor rewrites. Different
operations lower to scalar addition, multiplication, and summation, so laws
over these primitives can relate expressions built from different tensor
operations. TEPL normalizes both sides under the arithmetic laws permitted
by the selected policy.

Consider distributing a generalized dot over elementwise addition without
broadcasting, with identical dot descriptors on both sides:
\begin{equation}
 \operatorname{dot\_general}(X+Y,Z)
 \;\longleftrightarrow\;
 \operatorname{dot\_general}(X,Z)+\operatorname{dot\_general}(Y,Z).
\end{equation}
Addition's local value expression is \code{Add(X, Y)} under group $S$.
The dot's local expression is \code{Reduce(K, Product(lhs, Z))} under its own
batch, contraction, and free groups. Substituting the arithmetic expression
for \code{lhs} gives the expected shape of the computation, but does not
establish how positions in $S$ correspond to $K$ or retain their pairing with
\code{Z}'s contraction axes.

Symbolic axis mapping supplies this missing correspondence. The dot's
contraction attribute selects positions in addition's output. Since addition
preserves axis order, those positions identify corresponding axes in both
\code{X} and \code{Y}. The compiler records this relationship, and similarly
connects batch and free axes, even when group lengths are unknown. With fixed
output coordinates and indexed reads abbreviated, the two expressions become
\begin{equation}
 \sum_K(X+Y)Z
 \qquad\text{and}\qquad
 \sum_K XZ+\sum_K YZ.
\end{equation}
In the value language they are
\begin{codeblock}[H]
\begin{lstlisting}[language=TEPL]
Reduce(K, Product(Add(X, Y), Z))

Add(Reduce(K, Product(X, Z)),
    Reduce(K, Product(Y, Z)))
\end{lstlisting}
\end{codeblock}
Here, operand names abbreviate reads under aligned bindings, and $K$ is the
corresponding contraction group on both sides. Distribution of multiplication
and summation over addition brings the expressions to a common form when
the arithmetic policy permits those laws. Section~\ref{sec:axis-mapping}
explains how the required mappings are represented, composed, and refined.


\section{Mapping Derivations and Construction Algorithm}
\label{app:mapping}
Symbolic axis mapping connects groups defined by different operations. It
records where each group appears in a tensor, allowing composition of value
expressions while preserving axis order and operand pairings.


\subsection{Building and Composing Mapping Tables}
\label{extended:sec:mapping-tables}
Each record $(\mathsf{I},G,q)$ satisfies Equation~\eqref{extended:eq:mapping}. We display records
in tables whose columns denote tensors, whose first row gives the current
interface, and whose other rows give group positions. A dash means that the
example has no mapping for that group. Tables can include both registered
and derivable mappings.


Figure~\ref{fig:axis-mapping} first illustrates the principle with
\code{T = dot(X,Y);} \code{O = add(T,Z)}. The dot describes its output using $I$ and
$J$, while addition describes that same interface using $S$. Connecting
$S\leftrightarrow I\catseq J$ through $T$ aligns the coordinate descriptions
without changing the arithmetic expression or physical axis order. The
concrete matrix positions below the symbolic expressions are examples,
rather than fixed-rank restrictions on the method. Unlike the following
worked example, the figure places the dot inside the addition; both require
connecting local descriptions through a shared tensor.


We now work through the distribution example
\code{T = Add(X,Y);} \code{O = dot_general(T,Z)} from
Section~\ref{extended:sec:composition}. Let $\iota(G)=\operatorname{range}(|G|)$ denote
all positions in $G$. Ranks and position attributes remain symbolic.
\paragraph{1. Build addition's interface.}
Addition binds \code{X} and \code{Y} to one ordered group $S$ and yields $S$.
Its reads and result thus have these mappings:
\begin{center}
\begin{tabular}{@{}lccc@{}}\toprule
Interface/group & $X$ & $Y$ & $T$\\\midrule
Axis interface & $S$ & $S$ & $S$\\
$S$ & $\iota(S)$ & $\iota(S)$ & $\iota(S)$\\\bottomrule
\end{tabular}
\end{center}



\paragraph{2. Establish the dot's operand layouts.}
The dot introduces batch group $B$, contraction group $K$, and free groups
$I,J$. Let $b_\ell,b_r$ be its batch attributes, $k_\ell,k_r$ its contraction
attributes, and $S_Z$ the initial interface of \code{Z}. Free positions are
\begin{equation}
\begin{aligned}
 f_\ell&=\ocomp(\iota(S),b_\ell\catseq k_\ell),\\
 f_r&=\ocomp(\iota(S_Z),b_r\catseq k_r).
\end{aligned}
\end{equation}
Bindings require valid position lists and matching paired extents and must
cover each operand without overlap. Groups can be interleaved, so concatenating
$B,K,I$ does not necessarily reproduce the left operand's dimension order.
Instead, binding positions restore the original order:
\begin{equation}
\begin{aligned}
 \widehat S&=\gath(B\catseq K\catseq I,
                  \operatorname{inverse}(b_\ell\catseq k_\ell\catseq f_\ell)),\\
 \widehat S_Z&=\gath(B\catseq K\catseq J,
                  \operatorname{inverse}(b_r\catseq k_r\catseq f_r)).
\end{aligned}
\label{extended:eq:rebuilt-interfaces}
\end{equation}
Each concatenated position sequence is a permutation because the bindings
form a complete disjoint partition. Thus
\begin{equation}
\begin{aligned}
 \gath(\widehat S,b_\ell)&=B,&
 \gath(\widehat S,k_\ell)&=K,&
 \gath(\widehat S,f_\ell)&=I.
\end{aligned}
\end{equation}
The right interface satisfies analogous bindings. Registering these mappings
and partitions gives
\begin{center}
\begin{tabular}{@{}lcc@{}}\toprule
Interface/group & $T$ & $Z$\\\midrule
Axis interface & $\widehat S$ & $\widehat S_Z$\\
$B$ & $b_\ell$ & $b_r$\\
$K$ & $k_\ell$ & $k_r$\\
$I$ & $f_\ell$ & --\\
$J$ & -- & $f_r$\\\bottomrule
\end{tabular}
\end{center}
The interfaces describe the operands in dot coordinates; they do not change
physical axis order. An existing producer interface is reused when it already
satisfies the bindings.

\paragraph{3. Apply the dot's mappings to \code{X} and \code{Y}.}
Addition describes \code{T} using $S$, whereas the dot describes those same
axes using $\widehat S$. Since addition preserves order, each axis of its
output corresponds to the same-position axes in \code{X} and \code{Y}.
Transporting its reads into $\widehat S$ therefore places $B,K,I$ at positions
$b_\ell,k_\ell,f_\ell$ in all three tensors. The compiler can now compose
\code{Reduce(K, Product(Add(X,Y), Z))}, with $K$ referring to corresponding
coordinates in \code{X}, \code{Y}, and \code{Z}.

\paragraph{4. Extend the table to the output.}
A later consumer sees the dot's output, where surviving axes occur in output
order and $K$ is absent. Registering the order $B\catseq I\catseq J$ extends
the table as follows. The row $\widehat S$ denotes addition's whole-axis group
after coordinate transport.
\begin{center}
\small
\setlength{\tabcolsep}{4pt}
\begin{tabular}{@{}lccccc@{}}\toprule
Interface/group & $X$ & $Y$ & $T$ & $Z$ & $O$\\\midrule
Axis interface & $\widehat S$ & $\widehat S$ & $\widehat S$ & $\widehat S_Z$ & $B\catseq I\catseq J$\\
$\widehat S$ & $\iota(\widehat S)$ & $\iota(\widehat S)$ & $\iota(\widehat S)$ & -- & --\\
$B$ & $b_\ell$ & $b_\ell$ & $b_\ell$ & $b_r$ & $\iota(B)$\\
$K$ & $k_\ell$ & $k_\ell$ & $k_\ell$ & $k_r$ & --\\
$I$ & $f_\ell$ & $f_\ell$ & $f_\ell$ & -- & $|B|+\iota(I)$\\
$J$ & -- & -- & -- & $f_r$ & $|B|+|I|+\iota(J)$\\\bottomrule
\end{tabular}
\end{center}
Adding a length to a position sequence shifts each entry. The entries for
\code{X} and \code{Y} follow from transporting the dot partition into
addition's reads; output entries follow from the declared order.

\paragraph{5. Find positions through existing mappings.}
Not every relationship needs to be stored directly. Suppose $q_X$ locates
$\widehat S$ in \code{X}, and $k_\ell$ locates $K$ in $\widehat S$:
\begin{equation}
 \gath(\mathsf{I}_X,q_X)=\widehat S,\qquad \gath(\widehat S,k_\ell)=K.
\end{equation}
Composing the selections yields
\begin{equation}
\begin{aligned}
 K&=\gath(\widehat S,k_\ell)\\
  &=\gath(\gath(\mathsf{I}_X,q_X),k_\ell)\\
  &=\gath(\mathsf{I}_X,\gath(q_X,k_\ell)).
\end{aligned}
\label{extended:eq:gather-composition}
\end{equation}
Thus the $K,X$ entry is $\gath(q_X,k_\ell)$. For addition, $q_X$ lists all
positions in order, simplifying this expression to $k_\ell$. Repeated composition
recovers positions through longer chains. Addition's single group makes the
correspondence straightforward; selecting across several producer groups
requires partitioning them to expose finer relationships.


\subsection{Symbolic Axis Partitioning}
\label{extended:sec:partitioning}
Figure~\ref{fig:axis-partitioning} illustrates the information retained when a
consumer selects across producer groups. Three positions must be distinguished:
a selected axis's position within its producer group, its slot in the
consumer's contraction list, and its partner's position in the other operand.



\paragraph{1. Why reassociation requires splitting groups.}
Dot reassociation changes $(XY)Z$ to $X(YZ)$. For generalized dots it must also
specify the axes contracted by each replacement. Consider this rule, using
\code{dot} for a dialect's generalized dot:
\begin{codeblock}[H]
\begin{lstlisting}[language=TEPL]
rule associate {
  (dot[@outer] (dot[@inner] X Y) Z)
  => (dot[auto] X (dot[auto] Y Z))
  where {
    @outer.lhs_batching_dimensions == @inner->B
  }
}
\end{lstlisting}
\end{codeblock}
Incompatible batch routing can change a reassociation's result. The guard
requires the outer dot's left batch list to equal the positions of the inner
dot's batch group in its result, in the same order. The expression
\code{@inner->B} denotes those positions. This restricts applicability;
the completed replacement still requires the checks in
Sections~\ref{sec:synthesis} and~\ref{sec:correctness}.

Let $T=\operatorname{dot\_general}(X,Y)$. The inner dot contracts $K$ and
yields $\mathsf{I}_T=B\catseq I\catseq J$, with free groups $I$ from \code{X} and
$J$ from \code{Y}. Write $k_\ell^{\mathrm{outer}},k_r^{\mathrm{outer}}$ for the
outer contraction lists in \code{T} and \code{Z}. Their shared group is
\begin{equation}
 D=\gath(\mathsf{I}_T,k_\ell^{\mathrm{outer}}).
\end{equation}
Since $B$ remains batch, the left contraction selects only from $I,J$.
Unlike addition's single group, this producer has groups from different
operands. The outer dot can select part of each, mixing the selected axes
in $D$. To build \code{dot_general(Y,Z)}, the compiler must separate the
portion from $J$, which can be contracted there, from the partners of $I$,
which must survive for the final dot with \code{X}.

\paragraph{2. Find the selected part of each producer group.}
For each producer group, retain both within-group positions and contraction
slots. Group $I$ occupies
\begin{equation}
 q_I=|B|+\iota(I)
\end{equation}
in $\mathsf{I}_T$. Compare this list to the outer selection using
\begin{equation}
 (u_I,v_I)=\mpo(k_\ell^{\mathrm{outer}},q_I).
\end{equation}
For distinct-entry sequences $a,b$, the operation
$\mpo(a,b)=(u,v)$ scans $a$ in order. Whenever $a_i=b_j$, it appends $i$ to
$u$ and $j$ to $v$. It returns positions in the input lists, rather than
the shared values. Hence $\gath(a,u)=\gath(b,v)$, with both results empty
when there is no overlap. For example,
\begin{equation}
\begin{aligned}
 a&=\langle3,0,4\rangle,&b&=\langle0,2,3\rangle,\\
 u&=\langle0,1\rangle,&v&=\langle2,0\rangle.
\end{aligned}
\end{equation}
The two gathers produce $\langle3,0\rangle$, retaining $a$'s order.
For symbolic inputs, the results remain position expressions until they can
be evaluated or simplified.

Here $u_I$ locates the selected axes in the contraction list, and $v_I$
locates them within $I$:
\begin{equation}
 \gath(k_\ell^{\mathrm{outer}},u_I)=\gath(q_I,v_I).
\end{equation}
The selected subgroup is therefore
\begin{equation}
 \Isel=\gath(I,v_I)=\gath(D,u_I).
\end{equation}
The remaining view uses an ordered complement:
\begin{equation}
 w_I=\ocomp(\iota(I),v_I),\qquad \Irem=\gath(I,w_I).
\end{equation}
The selected view follows outer contraction order, whereas the remaining
view follows $I$'s original order. Their concatenation need not reproduce
$I$, so the position sequences must be retained. The same construction
applies to $J$ using $q_J=|B|+|I|+\iota(J)$. For either $G\in\{I,J\}$,
$u_G,v_G$ locate its selected portion, and
$w_G=\ocomp(\iota(G),v_G)$ locates its remaining portion. The lengths and
selected portions remain symbolic.

Partitioning adds the following relationships:
\begin{center}
\small
\begin{tabular}{@{}lcccc@{}}\toprule
Group & In $\mathsf{I}_T$ & In $I$ & In $J$ & In $D$\\\midrule
$\Isel$ & $\gath(q_I,v_I)$ & $v_I$ & -- & $u_I$\\
$\Irem$ & $\gath(q_I,w_I)$ & $w_I$ & -- & --\\
$\Jsel$ & $\gath(q_J,v_J)$ & -- & $v_J$ & $u_J$\\
$\Jrem$ & $\gath(q_J,w_J)$ & -- & $w_J$ & --\\\bottomrule
\end{tabular}
\end{center}
The $D$ column supplies the connection to \code{Z} through contraction slots.

\paragraph{3. Map selected axes to \code{Z}.}
The outer dot pairs its left and right lists entry by entry. Selecting slots
$u_I$ and $u_J$ from its right list therefore gives the positions in \code{Z}
paired with the selected subgroups:
\begin{equation}
 z_I=\gath(k_r^{\mathrm{outer}},u_I),\qquad
 z_J=\gath(k_r^{\mathrm{outer}},u_J).
\label{extended:eq:z-partners}
\end{equation}
In Figure~\ref{fig:axis-partitioning}, \code{Z} has axes $[z_0,f,z_1]$ and
right contraction positions $[2,0]$. The left list selects $[j_1,i_0]$.
Axis $j_1$ occupies contraction slot $0$ and pairs with $z_1$ at position $2$;
$i_0$ occupies slot $1$ and pairs with $z_0$ at position $0$.
The selected groups thus retain all three descriptions:
\begin{center}
\begin{tabular}{@{}lccc@{}}\toprule
Selected group & In producer group & Slots in $D$ & In $Z$\\\midrule
$\Isel$ & $v_I$ & $u_I$ & $z_I$\\
$\Jsel$ & $v_J$ & $u_J$ & $z_J$\\\bottomrule
\end{tabular}
\end{center}
Combining these entries with producer mappings describes the original
computation in terms of the refined groups without losing operand pairings.

\paragraph{4. Extend the original mapping table.}
The original graph is $T=\operatorname{dot\_general}(X,Y)$ followed by
$O=\operatorname{dot\_general}(T,Z)$. Write $b_X,k_X,i_X$ for the known
positions of $B,K,I$ in \code{X}, and $b_Y,k_Y,j_Y$ for those of $B,K,J$ in
\code{Y}. Composition locates each subgroup: for example,
$\gath(i_X,v_I)$ locates $\Isel$ in \code{X}. In \code{Z}, let $b_Z$ be the
batch positions and $f_Z$ the positions of free group $H$.

The inner reduction removes $K$ from $T$; the outer reduction removes
$\Isel,\Jsel$ from $O$. The surviving order is
\begin{equation}
 \mathsf{I}_O=B\catseq\Irem\catseq\Jrem\catseq H.
\label{extended:eq:source-output}
\end{equation}
Table~\ref{extended:tab:refined-mappings} combines the refined relationships. To keep
entries readable, define
\begin{equation}
\begin{aligned}
 i_s&=\gath(i_X,v_I),&i_r&=\gath(i_X,w_I),\\
 j_s&=\gath(j_Y,v_J),&j_r&=\gath(j_Y,w_J),\\
 t_{Is}&=\gath(q_I,v_I),&t_{Ir}&=\gath(q_I,w_I),\\
 t_{Js}&=\gath(q_J,v_J),&t_{Jr}&=\gath(q_J,w_J),\\
 o_I&=|B|+\iota(\Irem),\\
 o_J&=|B|+|\Irem|+\iota(\Jrem),\\
 o_H&=|B|+|\Irem|+|\Jrem|+\iota(H).
\end{aligned}
\end{equation}
These are abbreviations for the position expressions, not additional
attributes to be supplied by the user.
\begin{table}[htbp]
\caption{Refined mappings across the original reassociation graph. Each entry
gives positions in the column's interface. A dash indicates no mapping in
this example.}
\label{extended:tab:refined-mappings}
\centering
\begin{tabular}{@{}lccccc@{}}\toprule
Group/interface & $X$ & $Y$ & $T$ & $Z$ & $O$\\\midrule
Axis interface & $\mathsf{I}_X$ & $\mathsf{I}_Y$ & $\mathsf{I}_T$ & $\mathsf{I}_Z$ & $\mathsf{I}_O$\\
$B$ & $b_X$ & $b_Y$ & $\iota(B)$ & $b_Z$ & $\iota(B)$\\
$K$ & $k_X$ & $k_Y$ & -- & -- & --\\
$\Isel$ & $i_s$ & -- & $t_{Is}$ & $z_I$ & --\\
$\Irem$ & $i_r$ & -- & $t_{Ir}$ & -- & $o_I$\\
$\Jsel$ & -- & $j_s$ & $t_{Js}$ & $z_J$ & --\\
$\Jrem$ & -- & $j_r$ & $t_{Jr}$ & -- & $o_J$\\
$H$ & -- & -- & -- & $f_Z$ & $o_H$\\\bottomrule
\end{tabular}
\end{table}

The procedure generalizes to other producer--consumer pairs: compare position
sequences on the shared interface, identify overlap, and add mappings from
each overlapping portion into both parent groups. Each partition is internally
disjoint, while groups in different partitions can overlap.


\subsection{General Construction Algorithm}
\label{extended:sec:mapping-algorithm}
Algorithm~\ref{alg:mapping} in Section~\ref{sec:mapping-algorithm} processes producers before consumers, connects each
operation's bindings to its inputs, and extends mappings to the output.
Overlap introduces the finer subgroup relationships described above.
The input is a source graph and its operation definitions; the output comprises
tensor interfaces, composed value expressions, established mappings,
and partitions.

\paragraph{State.}
The state is $\Sigma=(V,\mathcal M,\Pi)$. For each tensor $t$,
$V[t]=(\mathsf{I}_t,e_t)$ stores its ordered interface and value expression.
$\mathcal M$ contains individual mapping entries $(\mathsf{I},G,q)$. The partition collection
$\Pi$ records which groups cover an interface completely and without overlap:
\begin{equation}
 \pi=\bigl(\mathsf{I},\langle(G_1,q_1),\ldots,(G_m,q_m)\rangle\bigr).
\end{equation}
An interface can have several partitions because different operations group
its axes differently. A partition changes the grouping used in reasoning;
it does not change tensor shape, rank, or axis order. $\mathcal M$ locates the groups,
whereas $\Pi$ says which groups jointly cover the interface exactly once.

For example, if addition produces $T$ with $S=\langle i,j\rangle$, then
$V[T]=(\mathsf{I}_T,\operatorname{Add}(X,Y))$ with $\mathsf{I}_T=S$,
$\mathcal M$ includes $(\mathsf{I}_T,S,\langle0,1\rangle)$, and $\Pi$ includes the single-member
partition $(\mathsf{I}_T,\langle(S,\langle0,1\rangle)\rangle)$. A consumer may split
$S$ into $I=\langle i\rangle$ and $J=\langle j\rangle$. Then $\mathcal M$ gains their
position mappings and $\Pi$ gains the two-member partition alongside the
original one. A shape $\langle d_i,d_j\rangle$ remains unchanged.

\paragraph{Invariants.}
Every mapping must select its group in the stated order:
\begin{equation}
 (\mathsf{I},G,q)\in \mathcal M\ \Longrightarrow\ \gath(\mathsf{I},q)=G.
 \label{extended:eq:mapping-invariant}
\end{equation}
Every partition member satisfies the same invariant, and
$q_1\catseq\cdots\catseq q_m$ must be a permutation of $\iota(\mathsf{I})$.
Thus each axis occurs exactly once within one partition, although different
partitions may overlap. In the addition example,
$\gath(\langle i,j\rangle,\langle0,1\rangle)=S$, and $\langle0,1\rangle$
covers every position. The sequence $\langle0,0\rangle$ would repeat one
axis and omit the other, violating the partition invariant.



\code{initialize_inputs} gives each graph input an interface and read
expression. \code{producer_first} ensures descriptions are available before
use. \code{resolve_bindings} evaluates symbolic position expressions and
establishes validity and pairing conditions. It reuses producer coordinates
where applicable and otherwise introduces local groups. Each operand's
bindings $(G_i,q_i)$ cover its axes.

\code{register_mappings} adds the individual group-position entries.
\code{build_interface} reuses $\mathsf{I}$ when the required groups already occupy the
required positions. Otherwise it constructs
\begin{equation}
 Q=\gath\bigl(G_1\catseq\cdots\catseq G_m,
        \operatorname{inverse}(q_1\catseq\cdots\catseq q_m)\bigr).
\end{equation}
The inverse arranges collected groups in tensor dimension order.
\code{refine_and_register} registers the new partition and compares it
with earlier partitions on the same interface. It uses \code{matchPositions}
to add overlap mappings. \code{transport} updates the producer's read
coordinates from $\mathsf{I}$ to $Q$ and carries mappings and partitions into that
description for the current consumer. Equal interfaces need no coordinate
update. Finally, \code{evaluate} and \code{output_positions} derive the
output interface and positions from the declared order, and \code{compose}
combines transported expressions using the operation's value definition.

\paragraph{Deriving mappings on demand.}
A missing mapping can be recovered by following established relationships:
\begin{equation}
\begin{aligned}
 (\mathsf{I},G,q)\in \mathcal M\ \land\ (G,H,h)\in \mathcal M
 \quad\Longrightarrow\quad
 \gath(\mathsf{I},\gath(q,h))=H.
\end{aligned}
\end{equation}
Lookup can return $\gath(q,h)$ or register the corresponding entry. Some
relationships propagate during construction and others are derived when
queried; enumerating every possible mapping in advance is unnecessary.
Local bindings and output order establish local mappings, paired position
matches establish overlaps, and Equation~\eqref{extended:eq:gather-composition}
composes them. These operations are designed to maintain the invariant;
they do not dispense with checking the completed rewrite.


\section{Synthesis Derivations and Algorithm}
\label{app:synthesis}
The mapping procedure establishes which logical axes occur in source tensors
and where they occur. Synthesis uses these relationships to fill the
replacement's attribute holes. It normalizes value expressions, identifies
which axes each replacement operation must preserve or reduce, and queries
mappings to construct its attributes.

Figure~\ref{fig:auto-matching} previews this process for distribution of a dot
over addition. The source is lowered into Axes IR and normalized; the RHS
arithmetic skeleton supplies reads to match before its attributes are known.
Read correspondence allows reuse of source coordinates, after which axis
positions construct the two missing dot descriptors. Here addition has no
broadcasting and preserves paired axis order, so both inferred descriptors
equal the source descriptor. The figure's reads $A[I,K]$, $B[I,K]$, and
$C[K,J]$ abbreviate group coordinates, rather than asserting that groups are
contiguous in physical layouts. The skeleton's question marks denote unresolved
reduction information. Equivalence is checked only after construction.



For attribute reconstruction beyond descriptor reuse, continue the guarded
reassociation example under an arithmetic policy permitting the required laws:
\begin{equation}
 \underbrace{\operatorname{dot}(\operatorname{dot}(X,Y),Z)}_{\mathrm{LHS}}
 \ \longrightarrow\
 \underbrace{\operatorname{dot}(X,\operatorname{dot}(Y,Z))}_{\mathrm{RHS}}.
\end{equation}
The rule supplies the replacement graph. Let $T=\operatorname{dot}(X,Y)$ be
the source intermediate, $U=\operatorname{dot}(Y,Z)$ the replacement
intermediate, and $O,O'$ their respective final results.


\subsection{Value Expression Normalization}
\label{extended:sec:normalization}
The graphs combine their reads in different orders. Normalization identifies
corresponding computations before constructing the replacement's attributes.
With output coordinates fixed and indexed reads abbreviated, the LHS is
\begin{equation}
 e_{\mathrm{LHS}}=\sum_D\left(\sum_K XY\right)Z.
\end{equation}
Each sum ranges over the coordinates of its group. The abbreviated notation
hides, but does not discard, the mappings established in
Section~\ref{sec:axis-mapping}.

TEPL uses a sum-of-products representation similar to SPORES's R-polynomial
form~\cite[Section~2.3, Definition~2.4]{spores}. A term contains a coefficient,
indexed operand reads, and symbolic reduction axes:
\begin{equation}
 t=(c,[r_1,\ldots,r_m],\mathcal D),\qquad
 \operatorname{value}(t)=\sum_{\mathcal D} c\prod_{i=1}^{m}r_i.
 \label{extended:eq:normalized-term}
\end{equation}
Each read retains its tensor identity and coordinate expression. $\mathcal D$ is an
ordered sequence of symbolic axes, and an empty sequence denotes no summation.
The empty product is $1$, permitting constants. A normalized expression is a
list $\mathcal E=[t_1,\ldots,t_n]$ denoting
\begin{equation}
 \operatorname{value}(\mathcal E)=\sum_{j=1}^{n}\operatorname{value}(t_j).
\end{equation}
Reads within a term multiply, separate terms add, and the empty term list
represents zero. Examples, again abbreviating indexed reads, are:
\begin{center}
\begin{tabular}{@{}ll@{}}\toprule
Expression & Normalized term list\\\midrule
$XY$ & \code{[(1, [X, Y], [])]}\\
$X(Y+Z)=XY+XZ$ & \code{[(1, [X, Y], []), (1, [X, Z], [])]}\\
$2XY-Z$ & \code{[(2, [X, Y], []), (-1, [Z], [])]}\\
$\sum_K XY$ & \code{[(1, [X, Y], K)]}\\\bottomrule
\end{tabular}
\end{center}
This similarity to SPORES concerns representation; it does not transfer a
completeness result to TEPL's supported reasoning.

Under the permitted multiplication and reduction laws, the example normalizes
to
\begin{equation}
 \sum_D\left(\sum_K XY\right)Z
 =\sum_D\sum_K XYZ
 =\sum_{D,K}XYZ.
\end{equation}
The inner reduction contributes reads \code{X}, \code{Y}, and axes $K$.
Multiplication incorporates \code{Z} after checking that it does not depend
on $K$. The outer reduction prepends $D$ after establishing that $D$ and $K$
are disjoint. Thus
\begin{codeblock}[H]
\begin{lstlisting}
Value expression:
Reduce(D,
  Product(
    Reduce(K, Product(X, Y)),
    Z))

Normalized term:
t = (1, [X, Y, Z], D ++ K)
\end{lstlisting}
\end{codeblock}

Unknown RHS attributes mean its actual reduction axes and read coordinates
are not yet available. Synthesis instead normalizes its \emph{arithmetic
skeleton}: it expands declared value operations while temporarily omitting
axis information. The skeleton for $X(YZ)$ has coefficient $1$, reads
\code{[X,Y,Z]}, and an empty reduction sequence. That empty sequence is a
property of the skeleton representation, not a claim that the completed RHS
has no reductions.

Synthesis matches each RHS read with the same operand's LHS read and reuses
its coordinates. For several terms, it first matches terms with the same
operands. Multiple possible read-coordinate assignments must be shown to
agree; otherwise synthesis remains unresolved. This correspondence provides
coordinates for construction. Reduction domains, coefficients, and full
value equivalence still have to be checked after attributes are known.


\subsection{Determine Which Axes Each Operation Preserves}
\label{extended:sec:keep-axes}
After read matching, the logical axes in each RHS input are known, but their
roles in the replacement may change. The source inner dot contracts $K$;
the source outer dot contracts $D$, partitioned into $\Isel$ and $\Jsel$.
The final output retains $B,\Irem,\Jrem,H$. Source mappings show that
\code{Y} contains $B,K,\Jsel,\Jrem$ and \code{Z} contains
$B,\Isel,\Jsel,H$. These lists specify membership rather than physical order.
To build $U$, the compiler must classify these axes as batch, contraction,
left-free, or right-free. In particular, $K$ must survive $U$ for a later
contraction with \code{X}.


Figure~\ref{fig:infer-u} shows the decision for an unbatched instance.
Starting with axes in \code{Y} or \code{Z}, it first preserves those needed
by \code{X}, then those required by the final output. The remaining shared
group, $\Jsel$, can be reduced inside $U$. The figure omits batch group $B$;
the equations below retain it for the guarded general case.

Apply the mandatory-axis definition in Equation~\eqref{eq:mandatory-axes}
and the live-set rules in Section~\ref{sec:canonical-axes} to this example.
Here $\Axes(x)$ is the set of logical axes in tensor $x$'s interface;
$\Axes_{\mathrm{in}}(U)$ and $\Axes_{\mathrm{out}}(U)$ collect axes used by
reads inside and outside $U$'s subtree.

For $U$, the inside set is $\Axes(Y)\cup\Axes(Z)$ and the outside set is
$\Axes(X)$. Using group names as shorthand for all their constituent axes,
\begin{equation}
\begin{aligned}
 \Axes_{\mathrm{in}}(U)
    &=\{B,K,\Isel,\Jsel,\Jrem,H\},\\
 \Axes_{\mathrm{out}}(U)
    &=\{B,K,\Isel,\Irem\},\\
 \Axes(O)&=\{B,\Irem,\Jrem,H\},\\
 M_{U}&=\{B,K,\Isel,\Jrem,H\}.
\end{aligned}
\end{equation}
Group $K$ occurs inside and outside and so is preserved. $\Jsel$ is needed
only inside and can be contracted. Although $\Irem$ belongs to the final
output, it occurs outside $U$ and reaches the final dot through \code{X}.
The resulting decisions are summarized in Table~\ref{tab:keep-decisions}.
At the final dot, $K$ and $\Isel$ are shared by \code{X} and $U$ but absent
from $O$, so they are reduced together.



\paragraph{From required axes to operation groups.}
The mandatory set identifies coordinates that must survive. Assigning coordinates to
operation groups additionally uses which operands contain them and whether
the declared value expression reduces the group. For an operation producing
$t$ from $x_1,\ldots,x_m$, let $\Axes_i=\Axes(x_i)$, and let $E_g$ be the
nonempty set of operands bound to declared group $g$. Axes appearing in exactly
those operands form
\begin{equation}
 D_g=\left(\bigcap_{i\in E_g}\Axes_i\right)
     \setminus\left(\bigcup_{i\notin E_g}\Axes_i\right).
\end{equation}
For roles that can reduce or preserve those coordinates, a candidate rule is
\begin{equation}
 \Axes_g=\begin{cases}
 D_g\setminus M_{t},&\text{if the value definition reduces }g,\\
 D_g\cap M_{t},&\text{otherwise.}
 \end{cases}
 \label{extended:eq:group-assignment}
\end{equation}
This candidate rule is subject to the full operation contract; it is not a
replacement for the canonical live-set rules. For a dot, shared axes in $M_t$
are batch and shared axes outside $M_t$ are contracted. All left-only and
right-only axes remain free, even if they are outside $M_t$. In this example,
the free axes are mandatory too, giving the following assignments:
\begin{center}
\small
\begin{tabular}{@{}lcccc@{}}\toprule
Operation & Batch & Contracted & Left free & Right free\\\midrule
$U=\operatorname{dot}(Y,Z)$ & $B$ & $\Jsel$ & $K,\Jrem$ & $\Isel,H$\\
$O'=\operatorname{dot}(X,U)$ & $B$ & $K,\Isel$ & $\Irem$ & $\Jrem,H$\\\bottomrule
\end{tabular}
\end{center}
Free axes follow operand order. The rule relies on distinguishable declared
group roles; ambiguous assignments remain unresolved. Groups introduced by an
explicit output domain rather than operand bindings are handled separately
in Section~\ref{extended:sec:synthesis-algorithm}.


\subsection{Construct Attributes from the Mapping Table}
\label{extended:sec:attribute-construction}
Group membership says which axes a dot uses, but its attributes need ordered
positions in actual operands. A query $\pos(\mathsf{I},G)$ returns $q$ satisfying
$\gath(\mathsf{I},q)=G$, using existing mappings and their compositions. Both sides
of a paired binding query the same \emph{ordered} group, ensuring that
corresponding entries select paired axes.


Figure~\ref{fig:attribute-construction} shows the two stages for the unbatched
example. The first pair of queries constructs the inner replacement's
contraction lists. Its completed interface $\mathsf{I}_U$ then supplies mappings for
the final dot's queries. In the figure, $G$ abbreviates
$\Isel\catseq K$, and preserved-group chips specify membership, not order.
The short fields \code{lhs_contract} and \code{rhs_contract} abbreviate the
contraction attributes. The following derivation also includes batch lists.


\paragraph{1. Construct the inner replacement dot.}
For $U=\operatorname{dot}(Y,Z)$, the source table already locates $B$ and
$\Jsel$ in both operands. Its attributes are
\begin{equation}
\begin{aligned}
 b_\ell^U&=b_Y,&b_r^U&=b_Z,\\
 k_\ell^U&=\gath(j_Y,v_J),&k_r^U&=z_J.
\end{aligned}
\end{equation}
By Equation~\eqref{extended:eq:z-partners}, $z_J=\gath(k_r^{\mathrm{outer}},u_J)$,
so both lists select $\Jsel$ in the same order. Free positions are
\begin{equation}
\begin{aligned}
 f_Y^U&=\\
 &\ocomp(\iota(\mathsf{I}_Y),b_Y\catseq k_\ell^U),\\
 f_Z^U&=\\
 &\ocomp(\iota(\mathsf{I}_Z),b_Z\catseq k_r^U),
\end{aligned}
\end{equation}
and the dot's output order constructs
\begin{equation}
\begin{aligned}
 \mathsf{I}_U&=B\catseq\gath(\mathsf{I}_Y,f_Y^U)\\
 &\catseq\gath(\mathsf{I}_Z,f_Z^U).
\end{aligned}
 \label{extended:eq:u-interface}
\end{equation}
The compiler extends mappings to $\mathsf{I}_U$ using this layout. This is necessary
because $K$ and $\Jrem$ can be interleaved in \code{Y}, as can $\Isel$ and
$H$ in \code{Z}. The interface preserves those physical orders instead of
assuming that each surviving group is contiguous.



\paragraph{2. Construct the final replacement dot.}
The final dot contracts $K$ and $\Isel$. Filtering source reduction order
$D\catseq K$ to those axes gives $\Isel\catseq K$: axes from $I$ occur
within $D$, before $K$. Use that same order for both contraction lists:
\begin{equation}
\begin{aligned}
 k_\ell^{O'}&=\gath(i_X,v_I)\catseq k_X,\\
 k_r^{O'}&=\pos(\mathsf{I}_U,\Isel)\catseq\pos(\mathsf{I}_U,K).
\end{aligned}
\end{equation}
A fixed binding determines group order, and inferred paired lists must follow
it. Otherwise the implementation prefers source reduction order for reduced
groups and source output order for retained groups, appending remaining axes
in operand order. Table~\ref{extended:tab:replacement-attributes} completes the two
descriptors.
\begin{table}[htbp]
\caption{Candidate attributes for the replacement dots.}
\label{extended:tab:replacement-attributes}
\begin{tabularx}{\linewidth}{@{}lYY@{}}\toprule
Attribute & $U=\operatorname{dot}(Y,Z)$ & $O'=\operatorname{dot}(X,U)$\\\midrule
Left batch & $b_Y$ & $b_X$\\
Right batch & $b_Z$ & $\iota(B)$\\
Left contraction & $\gath(j_Y,v_J)$ & $\gath(i_X,v_I)\catseq k_X$\\
Right contraction & $z_J$ & $\pos(\mathsf{I}_U,\Isel)$\newline $\catseq\pos(\mathsf{I}_U,K)$\\\bottomrule
\end{tabularx}
\end{table}
Removing the final contraction axes leaves
\begin{equation}
 \mathsf{I}_{O'}=B\catseq\Irem\catseq\Jrem\catseq H=\mathsf{I}_O.
\end{equation}
These attributes therefore have the candidate pairings and output order.
Operation legality and value equivalence still require checking.


\subsection{Produce and Verify Attribute Programs}
\label{extended:sec:emit-programs}
The tables contain symbolic group queries, while a rule application must
compute concrete attributes. The compiler resolves each query to expressions
over source attributes and shapes. For the inner replacement's right
contraction list, for example,
\begin{equation}
\begin{aligned}
 u_J&=\mpo\bigl(k_\ell^{\mathrm{outer}},
                   |B|+|I|+\iota(J)\bigr)_1,\\
 k_r^U&=\gath(k_r^{\mathrm{outer}},u_J),
\end{aligned}
\label{extended:eq:source-only-example}
\end{equation}
where subscript $1$ selects the first returned sequence. Group lengths and
ranges are themselves derived from source metadata. This expression constructs
positions at each match without accessing tensor elements. Queries involving
$\mathsf{I}_U$ similarly substitute its layout and compose mappings until only source
metadata remains. Emitted programs contain neither unresolved RHS attributes
nor logical-group queries.

For calculated accesses, position mappings alone need not determine every
attribute. The implementation additionally matches supported index expressions,
including coordinate coefficients and offsets. Composed slices, for example,
have the candidate attributes
\begin{equation}
\begin{aligned}
 \mathrm{start}'&=\mathrm{start}_{\mathrm{inner}}
   +\mathrm{stride}_{\mathrm{inner}}\,\mathrm{start}_{\mathrm{outer}},\\
 \mathrm{stride}'&=\mathrm{stride}_{\mathrm{inner}}\,
                  \mathrm{stride}_{\mathrm{outer}}.
\end{aligned}
\end{equation}
These expressions must also satisfy resulting access bounds. Coupled or
inverse equations outside the supported fragment remain unresolved.

\paragraph{Check the completed replacement.}
With attributes filled, the RHS has actual coordinates and reduction domains.
In the example, $U$ reduces $\Jsel$ and the final dot reduces
$\Isel\catseq K$:
\begin{equation}
\begin{aligned}
 e_{\mathrm{RHS}}
 &=\sum_{\Isel,K}X\left(\sum_{\Jsel}YZ\right)\\
 &=\sum_{\Isel,K,\Jsel}XYZ.
\end{aligned}
\end{equation}
The selected subgroups partition $D$, and the mappings preserve coordinate
pairing. The normalized RHS can therefore be compared to
$\sum_{D,K}XYZ$ under the permitted reduction laws. Verification also checks
output order, shape, dtype, and operation legality.

After construction, a fresh verification context reconstructs the completed
rule using the emitted programs. Required conditions must follow from legal
source metadata and the declared guard; construction obligations do not
become new assumptions. On success, the result is a collection of checked
source-dependent programs, compiled once and evaluated on each matched
graph's metadata. Unestablished construction or verification relationships
leave the rule unresolved.


\subsection{General Synthesis Algorithm}
\label{extended:sec:synthesis-algorithm}
Algorithm~\ref{alg:synthesis} in Section~\ref{sec:synthesis-algorithm} uses each operation's declared groups, value expression, and
output order, rather than requiring the replacement to contain dots. Its
inputs are a rule's source graph, replacement graph with attribute holes,
operation definitions, and declared conditions. It returns checked attribute
programs or an unresolved result.

The state retains $(V,\mathcal M,\Pi)$ and adds a table $\mathcal C$ of proposed attribute
expressions indexed by operation and attribute. In the pseudocode, \code{.axes} and \code{.value} denote the two
components of $V[t]$. Read correspondence supplies required output sets,
and replacement nodes are processed after their inputs. The construction
invariant requires every completed node to have resolved legal attributes,
an interface and value expression consistent with them, and mappings
satisfying Equation~\eqref{extended:eq:mapping-invariant}. Whole-rule equivalence is
checked after completion.



\code{correspond} matches normalized reads, and \code{required_output_axes}
computes $M_{t}$ for each replacement intermediate. \code{select_groups}
uses operand membership and reduction roles as in
Section~\ref{extended:sec:keep-axes}. Explicit output-domain groups are separate: they
come from required output axes absent from every operand, and their sizes
determine domain attributes. Ambiguous group assignments remain unresolved.
\code{order_groups} respects fixed bindings and otherwise follows source
and operand orders. \code{positions} implements $\operatorname{pos}$ by mapping lookup and position
construction for these ordered groups.

\code{complete_other_attributes} handles supported index coefficients,
offsets, output domains, and output permutations. \code{lower_and_check}
checks each completed operation and publishes its interface and mappings for
consumers. \code{check_equivalence} compares fully lowered values, reduction
domains, output order, shape, and dtype. \code{source_only_expressions}
resolves constructors to source metadata, and the final check reconstructs
the completed rule in a fresh context. The pseudocode omits failure paths:
unsupported or unproved steps return an unresolved result. Successful output
is a program for every missing attribute, ready for match-time evaluation.
